% Generated draft: requires mathematical and source review.
\documentclass[11pt]{article}
\usepackage[margin=25mm]{geometry}
\usepackage{fontspec}
\setmainfont{Latin Modern Roman}
\usepackage{amsmath,amssymb,mathtools,unicode-math}
\setmathfont{Latin Modern Math}
\usepackage{enumitem,needspace,xurl,hyperref}
\setlength{\parindent}{0pt}
\setlength{\parskip}{5pt}
\setlength{\emergencystretch}{3em}
\newcommand{\N}{\mathbb{N}}
\newcommand{\Z}{\mathbb{Z}}
\newcommand{\Q}{\mathbb{Q}}
\newcommand{\R}{\mathbb{R}}
\newcommand{\C}{\mathbb{C}}
\newcommand{\F}{\mathbb{F}}
\newcommand{\A}{\mathbb{A}}
\newcommand{\PP}{\mathbb{P}}
\newcommand{\E}{\mathbb{E}}
\newcommand{\eps}{\varepsilon}
\newcommand{\dd}{\,\mathrm{d}}
\begin{document}
\section*{UnsolvedMath ID: 30001192}
\textbf{Source problem code:} catalogue entry
\subsection{Definitions and mathematical statement}
The following formulates the question using the usual modular representation conventions. These conventions, including the meaning of a \(p\times p\) partition and the ground field, are not explicitly defined in the supplied record and require confirmation against the report.

Let \(p\) be a prime, let \(k\) be an algebraically closed field of characteristic \(p\), and let \(n\geq 1\). Write a partition of \(n\) as
\[
\lambda=(a_1^{m_1},\ldots,a_t^{m_t}),\qquad a_1>\cdots>a_t>0,
\]
where \(m_j\) is the multiplicity of the part \(a_j\). Under the standard convention used here, \(\lambda\) is a \(p\times p\) partition if
\[
p\mid a_j\quad\text{and}\quad p\mid m_j\qquad(1\leq j\leq t).
\]
Thus its Young diagram consists of \(p\times p\) squares arranged according to a partition; the term does not mean only the single square partition \((p^p)\).

Let \(\operatorname{core}_p(\lambda)\) be the \(p\)-core obtained by successively removing rim hooks of length \(p\). The block weight of \(\lambda\) is
\[
w=\frac{n-\lvert\operatorname{core}_p(\lambda)\rvert}{p}.
\]
Let \(S_k^\lambda\) denote the Specht module for the symmetric group \(\mathfrak S_n\), constructed over \(k\) from the partition \(\lambda\). No \(p\)-regularity assumption on \(\lambda\) is imposed.

For a finite-dimensional \(k\mathfrak S_n\)-module \(M\), let
\[
\cdots\longrightarrow P_2\longrightarrow P_1\longrightarrow P_0\longrightarrow M\longrightarrow0
\]
be a minimal projective resolution. Its complexity \(c_{k\mathfrak S_n}(M)\) is zero when \(M\) is projective; otherwise it is the least positive integer \(c\) for which some constant \(C>0\) satisfies
\[
\dim_k P_r\leq C r^{c-1}
\qquad\text{for all sufficiently large integers }r.
\]

\textbf{Question under these conventions.} For every prime \(p\), every such field \(k\), and every nonempty \(p\times p\) partition \(\lambda\) of positive block weight \(w\), must
\[
c_{k\mathfrak S_n}(S_k^\lambda)=w-1
\]
hold, or do there exist such partitions for which
\[
c_{k\mathfrak S_n}(S_k^\lambda)<w-1?
\]
Describe these cases further, in particular by characterizing any partitions exhibiting a strict drop and determining their complexities.

The nonempty, positive-weight scope is an explicit interpretation: the supplied question does not discuss the empty partition or weight zero. The source asks about equality versus a smaller value; this draft does not independently establish the upper bound implicit in that framing. The requested description of cases is not given a more specific classification target in the record.
\subsection{Short English statement}
For a \(p\times p\) partition \(\lambda\) of block weight \(w\), ask whether its Specht module always has complexity \(w-1\), or whether the complexity can be smaller. Describe any such cases more precisely. The modular setting and partition convention stated above require source confirmation; this is a statement-only draft.
\subsection{Sources}
\begin{itemize}
\item Local UnsolvedMath registry: permanent ID \(30001192\), title \emph{Complexity Drops for p-by-p Specht Modules}, exact source problem code \texttt{catalogue entry}; supplied complete source record in the Oberwolfach Reports Open Problems collection.
\item Supplied report attribution: \emph{Mini-Workshop: Support Varieties} (2009), Oberwolfach Reports, pp.\ 549--592, DOI \texttt{10.4171/owr/2009/10}. \url{https://doi.org/10.4171/owr/2009/10}. The record gives proposed year 2009 and names no proposer. It describes Oberwolfach Reports as documenting mathematical workshops and their open-problem discussions.
\item Supplied literature reference: Hemmer 2009. \url{https://doi.org/10.1007/s10801-009-0169-0}.
\item Supplied literature reference: Lim 2009. \url{https://doi.org/10.1016/j.jalgebra.2009.01.016}.
\item Supplied literature reference labelled ``original Oberwolfach problem'': \url{https://doi.org/10.4171/owr/2009/10}.
\item Supplied literature assessment, checked \textbf{2026-08-21}: status ``partially solved.'' It reports complexity \(p-1\) for the basic \(p\times p\) partition and in several special families, while describing the possibility of a lower complexity for arbitrary \(p\times p\) partitions of weight \(w\) as open. This is supplied dated curation derived from the OWR record, not an independently verified literature assessment; the special families are not specified here.
\item Supplied statement-verification metadata: status \texttt{corrected\_verified}. The record states that the mathematical content was checked against the official report and that notation, context, and question form were normalized without changing the source-level claim. This draft reports that verification claim without independently confirming it.
\end{itemize}
\end{document}
